\section{Proof dependencies}
The optical recurrence, normalization and logarithmic sensitivities are
those of the retained positive optical polynomial theorem (module 182).
The all-contact and original-guard derivatives come from the retained
endpoint field and margin theorems (174, 177), and the full angular
factor uses the uniform cubic action and angular-density results
(180--181, 194). Lemma~\ref{lem:v93-cancellation} is the new absolute
use of that recurrence; it does not improve the logarithmic bound.

The measure parametrization and its incident-component convention are
Theorem \texttt{thm:v85-segment-map} in module 196. The signed
one-dimensional calculation is consistent with the distributional
$BV$ convention of \cite{AFP}. Its full fixed-count integrability is
proved in Theorem~\ref{thm:v93-traces}, using only qualitative finite
geometry there, never a collision-uniform word count. The historical
critical angular residue in module 176 is not represented as a new
observation.

The positive-cutoff source bound is Proposition
\texttt{prop:v85-cutoff}, especially equation
\texttt{eq:v85-cutoff-uniform}, in module 196. Its finite-count
remainder explains the retained order of the critical limit. The
complete fixed-width bad-source and positive scalar remainder are
the retained physical source theorem and module 108. The original
exceptional union, mass bound and Borel-prefix representation are
those of 204, 209, 212 and 215. Theorem~\ref{thm:v93-balance} adds
the sign and trace information rather than changing that selector.
The hereditary spectral interfaces remain mathematical inputs, not
certificates supplied by source hashes or finite diagnostics.
